\subsection{Optimal Privacy Budget Allocation}
\label{ss:optimal_budget}

The server now needs to make a decision at time $t$ whether it should
approximate the current statistics using the last released statistics (referred
to as \emph{approximation}), or publish the current statistics with the necessary
noise (referred to as \emph{publication}).
Both options will introduce errors in the released statistics, i.e., the
approximation incurs error $\hdis_t$, and the publication incurs error
$\var(\epsilon_{t,2})$ which is the estimation error of running FO with privacy
budget $\epsilon_{t,2}$.

Existing methods~\cite{Kellaris2014,Ren2022} directly compare $\hdis_t$ with
$\var(\epsilon_{t,2})$, and if
$\hdis_t<\var(\epsilon_{t,2})$, then the server chooses approximation, otherwise it chooses
publication.
We notice that this commonly used strategy is myopic, focusing only on the current
timestamp and neglecting long-term overall accuracy of the released
statistics.
Recent empirical studies also show that these existing methods often perform
worse than even static baseline methods~\cite{Schaler2024}.
To address this weakness, we propose a novel strategy, i.e.,
\emph{\textbf{O}ptimal privacy \textbf{B}udget \textbf{A}llocation (OBA)}, that
aims to achieve high long-term accuracy of released statistics.

The idea of OBA is that, at every time $t$, we make a decision of either
approximation or publication by considering the benefit in a long-term period of
$w$ timestamps backwards, rather than only the current timestamp.
In more detail, as we already know the dissimilarities $\hdis_{t-w+1}, \ldots,
\hdis_t$, an optimal privacy budget allocation strategy in a time window of size
$w$ should choose publication on those timestamps with large dissimilarities
(say, top $k$ largest dissimilarities), as choosing approximation on these $k$
timestamps is only likely to introduce larger errors in the released statistics.
Hence, a good strategy should spend the remaining $\epsilon/2$ privacy budget on
these $k$ timestamps.
% , with each timestamp spending $\epsilon/(2k)$ privacy budget, to publish stream statistics with noise.
Once a timestamp is selected for publication, its approximation error $\hdis_t$ is replaced by a publication error. The goal then is to minimize the total publication error (the sum of variances) across these $k$ selected timestamps under a fixed total budget of $\epsilon/2$. Since the magnitude of the original $\hdis_t$ no longer affects that timestamp’s error and $\var(\cdot)$ is convex in the budget, the total publication error is minimized by dividing the budget equally among all the $k$ timestamps.
The remaining $w-k$ timestamps with small dissimilarities simply choose
approximation.
Finally, if the current time $t$ belongs to these $k$ timestamps, then the server
chooses publication with privacy budget $\epsilon_{t,2}=\epsilon/(2k)$;
otherwise, the server chooses approximation.
OBA is illustrated in \cref{fig:OBA}.

\begin{figure}[htp]
  \centering
  \begin{tikzpicture}[ycomb,font=\small,>=stealth, inner sep=1pt,
  bar/.style={line width=3pt},
  mybrace/.style={decorate,thick,decoration={calligraphic brace,raise=1pt}},
  txt/.style={font=\footnotesize},
  arr/.style={-Implies,double,double distance=2pt},
  ]

  % top left
  \draw[bar] plot coordinates{(0,.8) (.2,1) (.4,.4) (.6,.5)};
  \draw[bar,color=gray!50] plot coordinates{(-.4,.3) (-.2,.6)};
  \draw[bar,color=red] plot coordinates{(.8,.7)};

  \draw[<->] (-.6,1.1) node[above] {$\hdis_t$} -- (-.6,0) -- (1.2,0) node[right] {time};
  \node[anchor=north] at (.8,0) {$t$};

  \draw[mybrace] (0,1) -- (.8,1) node[midway,above=3pt] {$w$};

  \draw[arr] (1.7,.6) -- ++(.6,0) node[midway,above=2pt] {sort};

  % top right
  \begin{scope}[xshift=3.5cm]
    \draw[bar] plot coordinates{(0,1) (.2,.8) (.6,.5) (.8,.4)};
    \draw[bar,color=red] plot coordinates{(.4,.7)};
    \node[anchor=north] at (.4,0) {$t$};

    \draw (-.2,0) -- (1,0);

    \draw[densely dashed,blue] (.5,-.1) -- ++(0,1.3) node[above] {$k^*$};

    \draw[Bar->] (.48,1.1) -- ++(-.8,0) node[above,txt] {publication};
    \draw[Bar->] (.52,1.1) -- ++(.8,0) node[above,txt] {approximation};

    \draw[arr] (1.4,.6) -- ++(.5,0) node[right=2pt] {publication};
  \end{scope}

  % bottom left
  \begin{scope}[yshift=-2cm]
    \draw[bar] plot coordinates{(0,.8) (.2,1) (.4,.4) (.6,.5)};
    \draw[bar,color=gray!50] plot coordinates{(-.4,.3) (-.2,.6)};
    \draw[bar,color=red] plot coordinates{(.8,.2)};

    \draw[<->] (-.6,1.1) node[above] {$\hdis_t$} -- (-.6,0) -- (1.2,0) node[right] {time};

    \draw[mybrace] (0,1) -- (.8,1) node[midway,above=3pt] {$w$};

    \node[anchor=north] at (.8,0) {$t$};

    \draw[arr] (1.7,.6) -- ++(.6,0) node[midway,above=2pt] {sort};
  \end{scope}

  % bottom right
  \begin{scope}[xshift=3.5cm,yshift=-2cm]
    \draw[bar] plot coordinates{(0,1) (.2,.8) (.4,.5) (.6,.4)};
    \draw[bar,color=red] plot coordinates{(.8,.2)};
    \node[anchor=north] at (.8,0) {$t$};

    \draw (-.2,0) -- (1,0);

    \draw[densely dashed,blue] (.5,-.1) -- ++(0,1.3) node[above] {$k^*$};

    \draw[Bar->] (.48,1.1) -- ++(-.8,0) node[above,txt] {publication};
    \draw[Bar->] (.52,1.1) -- ++(.8,0) node[above,txt] {approximation};

    \draw[arr] (1.4,.6) -- ++(.5,0) node[right=2pt] {approximation};
  \end{scope}

\end{tikzpicture}

  \caption{Illustration of OBA}
  \label{fig:OBA}
  \Description{The figure demonstrates the Optimal Budget Allocation (OBA) strategy using two scenarios. 
The top row illustrates a "publication" case: the estimated dissimilarity at the current timestamp $t$ (highlighted in red) is high relative to the sliding window w. When the dissimilarities in the window are sorted in descending order, the value at $t$ falls within the top $k^*$ threshold, triggering a data publication. 
The bottom row illustrates an "approximation" case: the dissimilarity at $t$ is low. After sorting, the red bar falls outside the top $k^*$ range into the approximation region, meaning the current data will be approximated using historical values to save privacy budget.}
\end{figure}


The last problem is how to find an optimal $k$.
Let $\hdis_{(1)},\dots,\hdis_{(w)}$ denote the $w$ dissimilarities
$\hdis_{t-w+1}, \dots, \hdis_t$ sorted in descending order.
The \emph{cumulative error} of the released statistics in a time window of size
$w$ consists of two parts: (i) cumulative publication error at these $k$
timestamps due to running $k$ FOs each using a privacy budget
$\frac{\epsilon}{2k}$, thus incurring an error
$k\cdot\var(\frac{\epsilon}{2k})$, and (ii) the cumulative approximation error
on the remaining $w-k$ timestamps, incurring another error
$\sum_{i=k+1}^w\hdis_{(i)}$.
Furthermore, we should allow $k=0$, and in this case, the server always chooses
approximation in the current window, thus incurring an error
$\sum_{i=1}^w\hdis_{(i)}$.

Let $E_t(k)$ denote the cumulative error of the released statistics in the most
recent window of size $w$ at time $t$, expressing $k$ as a parameter.
Then, we have
\[
  E_t(k) =
  \begin{cases}
    \sum_{i=1}^w \hdis_{(i)}, & k=0, \\
    k \cdot \var(\frac{\epsilon}{2k}) + \sum_{i=k+1}^w\hdis_{(i)}, & 0<k\le w.
  \end{cases}
\]
Our goal is to find an optimal $k^*$ by minimizing $E_t(k)$, i.e.,
\begin{equation}\label{eq:optimization}
  k^* \in \arg\min_{0\le k\le w} E_t(k).
\end{equation}
As $k$ is in a finite set $\{0,\ldots,w\}$, a simple method to solve
Problem~\eqref{eq:optimization} is to enumerate each $k$ and choose one
minimizing $E_t(k)$.
The pseudo-code of OBA is given in \cref{alg:oba}.

\begin{algorithm}[htp]
  \caption{\textbf{O}ptimal Privacy \textbf{B}udget \textbf{A}llocation
    (OBA)\label{alg:oba}}
  \KwIn{Dissimilarities $\hdis_{t-w+1}, \ldots, \hdis_t$, privacy budget
    $\epsilon$, remaining privacy budget for current window
    $\epsilon_t^\mathit{rm}$} \KwOut{Optimal privacy budget $\epsilon_{t,2}$ for
    timestamp $t$}

  $[\hdis_{(1)},\ldots,\hdis_{(w)}]\gets
  \text{Sort}([\hdis_{t-w+1},\ldots,\hdis_t])$\tcp*{descending}\label{ln:sort}
  $[t_{(1)},\ldots,t_{(w)}]\gets$ timestamps of $[\hdis_{(1)},\ldots,\hdis_{(w)}]$\;
  $E_t[0]\gets\sum_{i=1}^w\hdis_{(i)}$\;\label{ln:sum1}
  \lFor{$k\gets 1$ \KwTo $w$}{%
    $E_t[k]\gets k\cdot\var(\frac{\epsilon}{2k})
    + \sum_{i=k+1}^w\hdis_{(i)}$}\label{ln:sum2}
  $k^*\gets\arg\min_k E_t[k]$\;
  \uIf(\tcp*[f]{if publication}){$t\in\{t_{(1)},\ldots,t_{(k^*)}\}$}{\label{ln:t}
    $\epsilon_{t,2}\gets\min(\epsilon_t^{\mathit{rm}},
    \frac{\epsilon}{2k^*})$\tcp*{publication using budget
      $\epsilon_{t,2}$}\label{ln:epsilon}
  }
  \lElse(\tcp*[f]{approximation using previous time step}){%
    $\epsilon_{t,2}\gets 0$}\label{ln:zero}
  \Return $\epsilon_{t,2}$\;
\end{algorithm}


It is noteworthy to mention that the ultimate goal of OBA is to decide whether
we choose approximation or publication at time $t$, and if it is publication,
how much privacy budget to spend.
In Line~\ref{ln:t}, if current timestamp $t$ belongs to the selected $k^*$
timestamps, then we choose publication using privacy budget $\epsilon_{t,2}$;
otherwise we choose approximation with no privacy budget (Line~\ref{ln:zero}).
In addition, $\epsilon_{t,2}$ should be upper bounded by the remaining privacy
budget $\epsilon_t^\mathit{rm}$ available in current window
(Line~\ref{ln:epsilon}).

\paragraph{Improving Efficiency}
A straightforward implementation of OBA requires sorting dissimilarities and enumerating all possible $k$, resulting in a per-timestamp time complexity of $O(w\log w + w^2) =
O(w^2)$.
When $w$ is large, such brute-force enumeration becomes computationally
expensive in streaming settings.
We present several tricks to improve its computational efficiency.
\begin{itemize}
\item{Incremental Maintenance.} 
 Instead of sorting dissimilarities from scratch (Line~\ref{ln:sort}) at each timestamp, we maintain the sorted sequence $\{\hdis_{(1)}, \dots, \hdis_{(w)}\}$ incrementally.
 By using a priority queue, the insertion of the new $\hdis_t$ and the deletion of the expired $\hdis_{t-w}$ can both be performed in $O(\log w)$ time.
\item{Recursive Derivation.} The computation of the cumulative error $E_t(k)$ (Line~\ref{ln:sum2}) can be optimized by reusing intermediate results.
The publication error term $k\cdot\var(\frac{\epsilon}{2k})$ is time
  invariant which depends only on $k$, thus it can be computed offline and stored.
  The publication error term can be computed incrementally while sorting (and similarly for the sum in Line~\ref{ln:sum1}).
  So calculating $E_t$ actually has time complexity $O(w)$.
\item{Early Termination.}
    The function $E_t(k)$ typically exhibits a unimodal structure (initially decreasing and then increasing) or monotonicity.
    Specifically, $E_t(k)$ decreases when the benefit of removing a large approximation error $\hdis_{(k+1)}$ outweighs the cost of added noise, and begins to increase when the remaining dissimilarities are small.
    This structure enables an early-stop strategy: we iterate $k$ starting from $0$ and terminate immediately once $E_t(k+1) > E_t(k)$. Let $k_{\mathit{stop}}$ denote the number of values evaluated before termination.
    Thus, we only need to evaluate $k_{\mathit{stop}}$ candidates instead of all $w$ options.
\end{itemize}
By combining these techniques, the overall time complexity is reduced to $O(\log w + k_{\mathit{stop}})$, ensuring scalability even with large window sizes.